%% kapa: exam=1 round=35 year=2024 period=2 form=A subject=accounting subject_ko=회계학 seq=3 number=43
%% source: Q-Net 감정평가사 기출문제 (cst00309) artlSeq=5213229 file=raw/1008/5213229/2교시 문제지 원본.pdf
%% conversion: kordoc 4.24.51 -> tools/gen1.py (auto); review=proofread-sonnet; flags=tab
\begin{question}{43}
(주)감평의 20×1년 기말재고자산에 대한 자료가 다음과 같다.
\par{\small\centering\begin{tabular}{|c|r|r|r|r|}\hline
항목 & \multicolumn{1}{c|}{원가} & \multicolumn{1}{c|}{확정판매계약가격} & \multicolumn{1}{c|}{일반판매가격} & \multicolumn{1}{c|}{현행대체원가} \tabularnewline\hline
제품 A & ￦1,000 & ￦900 & ￦950 & - \tabularnewline\hline
제품 B & 1,200 & - & 1,250 & - \tabularnewline\hline
원재료 A & 1,100 & - & - & ￦1,000 \tabularnewline\hline
원재료 B & 1,000 & - & - & 900 \tabularnewline\hline
\end{tabular}\par}
\begin{framedtext}
○ 제품 A는 모두 확정판매계약을 이행하기 위하여 보유하고 있으며, 제품 A와 제품 B는 판매시 계약가격 또는 일반판매가격의 10\%에 해당하는 판매비용이 소요될 것으로 예상된다. \par
○ 원재료 A를 이용하여 생산하는 제품은 원가 이상으로 판매될 것으로 예상된다. \par
○ 원재료 B를 이용하여 생산하는 제품의 원가는 순실현가능가치를 초과할 것으로 예상된다.
\end{framedtext}
모든 재고자산에 대해 항목별기준을 적용할 때 20×1년도에 인식할 재고자산평가손실은? (단, 재고자산 감모는 발생하지 않았으며, 기초재고자산평가충당금은 없다.)
\begin{choices}
  \item ￦300
  \item ￦335
  \item ￦350
  \item ￦365
  \item ￦380
\end{choices}
\end{question}
\answer{4}
